\documentclass[11pt]{beamer}

\setbeamertemplate{navigation symbols}{}
\setbeamertemplate{headline}{}
\setbeamertemplate{footline}{}
\setbeamertemplate{sidebar left}{}
\setbeamertemplate{sidebar right}{}
\setbeamertemplate{background canvas}[transparent]

\usecolortheme{default}
\usefonttheme{default}


\usepackage{fontspec}
\usepackage{unicode-math}
\setmainfont{Tex Gyre Pagella}
\setmathfont{Tex Gyre Pagella Math}

\usepackage{ifthen}
\usepackage{amsmath, amssymb}
\usepackage{tikz} 
\usepackage{pgfplots}
\pgfplotsset{compat=1.18}
\usepackage[d]{esvect}
\usetikzlibrary{shapes.misc}
\usetikzlibrary{arrows,arrows.meta}
\usetikzlibrary{calc,intersections, patterns, math}

\usetikzlibrary{angles,quotes,babel,shapes}
\usetikzlibrary{decorations}
\usetikzlibrary{decorations.pathmorphing}
\usetikzlibrary{arrows}
\usepgfplotslibrary{fillbetween}
\usetikzlibrary{animations}
\usetikzlibrary{decorations.text}
\usetikzlibrary{math}
\usetikzlibrary{3d}
\usetikzlibrary{trees}

\definecolor{pfeil}{RGB}{168,167,167}
\definecolor{petrol}{RGB}{0, 118, 136}
\definecolor{darkgoldenrod}{RGB}{184, 134, 11}
\colorlet{petrol-lighter}{petrol!40}
\colorlet{darkgoldenrod-lighter}{darkgoldenrod!40}

\setbeamercolor{background canvas}{bg=transparent}
\setbeamercolor{normal text}{fg=pfeil,bg=transparent}

\newcommand{\vektor}[1]{
\begin{pmatrix}
	#1
\end{pmatrix}
}

\tikzset{%
	add/.style args={#1 and #2}{to path={($(\tikztostart)!-#1!(\tikztotarget)$)--($(\tikztotarget)!-#2!(\tikztostart)$)\tikztonodes}}
}

\begin{document}

\begin{frame}[plain]
\tikzstyle{every picture}+=[remember picture]

\begin{align*} 
\underset{\begin{matrix} \tikz[baseline]\node (x11) {$x_1$}; \\ \tikz[baseline]\node (y11) {$y_1$}; \end{matrix}}%
	{\vektor{\tikz[baseline]\node (x1) {$x_1$}; \\ \tikz[baseline]\node (y1) {$y_1$}; \\ \tikz[baseline]\node (z1) {$z_1$};}} 
\times \underset{\begin{matrix} \tikz[baseline]\node (x22) {$x_2$}; \\ \tikz[baseline]\node (y22) {$y_2$}; \end{matrix}}%
	{\vektor{\tikz[baseline]\node (x2) {$x_2$}; \\ \tikz[baseline]\node (y2) {$y_2$}; \\ \tikz[baseline]\node (z2) {$z_2$};}}
&= \vektor{\color{green!70!black}y_1\cdot z_2 \color{black} - \color{blue} y_2\cdot z_1 \\[0.5em] 
	\color{cyan}z_1\cdot x_2 \color{black} - \color{orange} z_2\cdot x_1 \\[0.5em] 
	\color{magenta}x_1\cdot y_2 \color{black} - \color{teal} x_2\cdot y_1} \\
\underset{\begin{matrix} \tikz[baseline]\node (x33) {$2$}; \\ \tikz[baseline]\node (y33) {$5$}; \end{matrix}}
	{\vektor{\tikz[baseline]\node (x3) {$2$}; \\ \tikz[baseline]\node (y3) {$5$}; \\ \tikz[baseline]\node (z3) {$1$};}} 
\times \underset{\begin{matrix} \tikz[baseline]\node (x44) {$3$}; \\ \tikz[baseline]\node (y44) {$6$}; \end{matrix}}
	{\vektor{\tikz[baseline]\node (x4) {$3$}; \\ \tikz[baseline]\node (y4) {$6$}; \\ \tikz[baseline]\node (z4) {$2$};}}
&= \vektor{\color{green!70!black}5\cdot 2 \color{black} - \color{blue} 6\cdot 1 \\[0.5em] 
		\color{cyan}1\cdot 3 \color{black} - \color{orange} 2\cdot 2 \\[0.5em] 
		\color{magenta}2\cdot 6 \color{black} - \color{teal} 3\cdot 5}
= \vektor{4 \\[0.5em] -1 \\[0.5em] -3}
\end{align*}
\begin{tikzpicture}[overlay,thick, add=0.05 and 0.05, line width=16pt, line cap=round, opacity=0.4,pfeil]
	\draw [green!70!black] (y1) to (z2);
	\draw [blue] (y2) to (z1);
	\draw [cyan] (z1) to (x22);
	\draw [orange] (z2) to (x11);
	\draw [magenta] (x11) to (y22);
	\draw [teal] (x22) to (y11); 
	
	\draw [green!70!black] (y3) to (z4);
	\draw [blue] (y4) to (z3);
	\draw [cyan] (z3) to (x44);
	\draw [orange] (z4) to (x33);
	\draw [magenta] (x33) to (y44);
	\draw [teal] (x44) to (y33);
\end{tikzpicture}
\end{frame}
\end{document}
